\section{Appendix: Datasets}
\label{sec:datasets}

We evaluate \textsc{VOILA} on five datasets spanning standard VQA benchmarks, text-centric reasoning, cyber--physical perception, and long-horizon agentic memory.
Table~\ref{tab:dataset_stats} summarizes the dataset statistics and experimental splits used throughout this work.

\begin{table}[h]
\centering
\small
\caption{\textbf{Dataset statistics and experimental splits.} All splits are the standard published splits unless otherwise noted. LoCoMo-VQA is our GPT-4o-generated benchmark.}
\label{tab:dataset_stats}
\begin{tabular}{lrrr}
\toprule
\textbf{Dataset} & \textbf{Images} & \textbf{Questions} & \textbf{Split used} \\
\midrule
VQA-v2   & 81{,}434 & 214{,}354 & val \\
GQA      & 22{,}669 & 12{,}578  & balanced test \\
TextVQA  & 5{,}734  & 5{,}734   & val \\
FloodNet & 2{,}343  & 4{,}511   & test \\
LoCoMo-VQA & 700   & 3{,}500   & generated \\
\bottomrule
\end{tabular}
\end{table}
Together, these datasets allow us to assess accuracy--cost tradeoffs under diverse visual content, linguistic structure, and deployment constraints.

\subsection{VQA-v2, GQA, and TextVQA}
\label{sec:vqa_datasets}
\paragraph{VQA-v2.}
VQA-v2 \cite{goyal2017vqav2} is a large-scale visual question answering benchmark consisting of real-world images paired with free-form questions and short answers.
It covers a broad range of question types, including object recognition, counting, attributes, and commonsense reasoning, and is designed to reduce language priors through balanced answer distributions.
We use the standard train/validation splits.

\paragraph{GQA.}
GQA \cite{hudson2019gqa} focuses on compositional and relational reasoning over images.
Questions are generated from structured scene graphs and emphasize multi-step reasoning, spatial relations, and logical composition.
We use GQA to test whether VOI generalizes beyond surface-level recognition to structured visual reasoning.

\paragraph{TextVQA.}
TextVQA \cite{singh2019textvqa} evaluates text-centric visual understanding, requiring models to read and reason over scene text embedded in images.
This dataset stresses OCR-dependent perception and complements VQA-v2 and GQA by emphasizing linguistic grounding in visual text.

\subsection{FloodNet}
\label{sec:floodnet}
FloodNet \cite{rahnemoonfar2021floodnet} is an aerial imagery dataset captured by UAVs during Hurricane Harvey, designed for disaster damage assessment.
It includes a native VQA component in which free-form questions are paired with short answers about flood conditions visible in each image (e.g., \emph{``Is the road flooded?''}, \emph{``What is the condition of the building?''}, \emph{``How many flooded buildings are visible?''}).
Question--answer pairs are drawn directly from the original FloodNet release; no additional generation was performed.
We use the FloodNet test split and evaluate models on exact-match accuracy against the provided reference answers.
FloodNet serves as a proxy for cost-sensitive perception in bandwidth-constrained IoT and emergency-response deployments, where transmission of high-resolution UAV imagery is expensive and decisions must be made under latency constraints.

\subsection{LoCoMo Question Generation.}
\label{sec:locomo}
LoCoMo~\cite{wang2024locomo} evaluates long-horizon agentic memory retrieval, where an agent must decide when expensive visual memory recall is necessary.
Unlike static VQA datasets, LoCoMo natively provides only (i) caption-level summaries cached in memory and (ii) the original image available via costly on-demand recall.

To enable quantitative evaluation, we construct a VQA-style benchmark over LoCoMo memory items.
Starting from LoCoMo conversational memory entries containing an image URL and associated evidence text, we generate a total of 3,500 question--answer pairs.
For each image, we produce multiple questions that are explicitly constrained to:
(i) reference visual content present in the image,
(ii) be relevant to the conversational evidence text, and
(iii) require varying degrees of visual detail to answer.

Question generation is performed using GPT-4o conditioned jointly on the image and its evidence text.
For each image, we generate a diverse set of questions spanning object counting, yes/no verification, spatial relations, attribute identification, and scene-level understanding using the following prompt:

\begin{quote}
\small
\texttt{Based on this image and the provided evidence text, generate exactly 5 diverse Visual Question Answering (VQA) questions with their answers.}

\texttt{Evidence text: \{evidence\_text\}}

\texttt{Generate questions in these categories:}\\
\texttt{1. Object counting (e.g., "How many people are in the image?")}\\
\texttt{2. Yes/No question (e.g., "Is there a dog in the image?")}\\
\texttt{3. Spatial awareness (e.g., "Where is the person standing?")}\\
\texttt{4. Object identification (e.g., "What color is the shirt?")}\\
\texttt{5. Scene understanding (e.g., "What activity is happening?")}

\texttt{Return ONLY a JSON array with 5 objects, each having "question" and "answer" fields. Example format:}\\
\texttt{[}\\
\texttt{\ \ \{\{"question": "How many people are visible?", "answer": "2"\}\},}\\
\texttt{\ \ \{\{"question": "Is there a tree in the background?", "answer": "yes"\}\},}\\
\texttt{\ \ ...}\\
\texttt{]}
\end{quote}

All generated questions are paired with answers produced by GPT-4o and filtered by the authors to remove questions that (i) cannot be answered from the image alone, (ii) contain hallucinated object references, or (iii) have ambiguous or non-unique answers.
Of the 4{,}375 initially generated QA pairs (875 images × 5 questions), 875 (20\%) were discarded during filtering, yielding 3{,}500 final pairs.
We note two important caveats of this construction.
First, using GPT-4o as both the question generator and the reference-answer provider means ground-truth answers reflect GPT-4o's visual interpretation; hallucinated answers propagate into the evaluation metric.
Second, the prompt enforces exactly one question per category per image, producing a uniform question-type distribution (20\% each) that differs from natural VQA distributions where yes/no questions are overrepresented.
We use LoCoMo exclusively as an out-of-distribution stress test and not as a primary benchmark; these construction choices are further discussed in the Limitations section.

This construction yields a LoCoMo-derived VQA benchmark that preserves the core constraint of agentic memory systems---cheap semantic summaries with expensive high-fidelity recall---while enabling controlled measurement of adaptive fidelity selection.
All LoCoMo experiments are conducted in a fully out-of-distribution setting: \textsc{VOILA} is trained on VQA-v2 and GQA and evaluated on LoCoMo without any task-specific retraining.
